\subsubsection{Subproceso 1: Acreditación y Pre-requisitos}

\textbf{Fuentes documentales:}
\begin{itemize}
    \item \textbf{TUPA de la UNI}~\cite[págs.~15--16]{tupa2024}: Establece que, antes de matricularse en la Intranet DIRCE, el alumno debe pagar el autoseguro médico y no tener adeudos en su Facultad (ver Figura~\ref{fig:ev1_tupa}).
    \item \textbf{Reglamento del Fondo de Autoseguro Estudiantil}~\cite[nums.~6.1.1 y 6.2]{autoseguro2023}: Regula el aporte estudiantil obligatorio que se abona en cada matrícula semestral (ver Figura~\ref{fig:ev1_autoseguro}).
    \item \textbf{Reglamento de Matrícula de Pregrado}~\cite[Caps.~I--V]{reglamento2022}: Define las reglas de cierre previo de notas, tutoría obligatoria, orden de mérito, turnos de matrícula, cambios de plan y requisitos de habilitación (ver Figura~\ref{fig:ev1_reglamento}).
    \item \textbf{Ley Universitaria y Estatuto de la UNI}~\cite{ley30220}: Base legal general del procedimiento de matrícula regular.
\end{itemize}

\textbf{Fuentes tecnológicas:}
\begin{itemize}
    \item \textbf{Portal Intralu (\url{https://intralu.uni.edu.pe/}):} Revisado para verificar la generación de órdenes de pago del autoseguro estudiantil y el registro de deudas del alumno (ver Figura~\ref{fig:ev1_intralu}).
    \item \textbf{Intranet DIRCE (\url{https://dirce.uni.edu.pe/}) y Sistema SIGA-ORCE:} Observados para comprobar cómo se muestran la ficha de matrícula, la tabla de grupos de turno de la Facultad y los controles de bloqueos, adeudos y pagos (ver Figuras~\ref{fig:ev1_requisitos} y~\ref{fig:ev1_ficha}).
\end{itemize}

\textbf{Fuentes humanas:}
\begin{itemize}
    \item \textbf{Estudiantes regulares de la Facultad de Ciencias (periodo 2026-2):} Alumnos de las escuelas de Ciencia de la Computación, Física e Ingeniería Física que reportaron cómo funcionó la habilitación de sus fichas, la validación de pagos y deudas, la convalidación por cambio de plan y la asignación de turnos en el ciclo 2026-2.
\end{itemize}
